\documentclass[12pt]{article}
\usepackage[utf8]{inputenc}
\usepackage[T1]{fontenc}
\usepackage[margin=1in]{geometry}
\usepackage{setspace}
\usepackage{xcolor}
\usepackage[colorlinks=true, urlcolor=blue!60!black, linkcolor=blue!60!black]{hyperref}
\setlength{\parindent}{2em}
\setlength{\parskip}{1.25em}
\renewcommand{\baselinestretch}{1.0}
\begin{document}
\begin{center}
{\LARGE \textit{National Republican Senatorial Committee v. Federal Election Commission}, 609 U.S. \_\_\_ (2026)\par}
\vspace{0.8em}
{\large Case Brief\par}
\end{center}
\vspace{1em}

\section{Memory Jogger}
In a 6--3 decision overruling \textit{Colorado II}, the Supreme Court held that the Federal Election Campaign Act's limits on political parties' coordinated expenditures with their own candidates violate the First Amendment because base contribution limits, earmarking rules, and modern disclosure requirements already address \textit{quid pro quo} corruption without imposing a disproportionate ``fourth prophylaxis'' on core political speech.

\section{Detailed Case Facts}
The Federal Election Campaign Act of 1971 (FECA), as amended in 1974 and codified at 52 U.S.C. \S~30116(d), caps the amount of money a political party committee may spend in ``coordination'' with its own federal candidates (\href{https://www.supremecourt.gov/opinions/25pdf/24-621_h315.pdf}{Slip Opinion, pp. 1--3}). Under Federal Election Commission (FEC) regulations, a coordinated expenditure is money spent ``in cooperation, consultation or concert with, or at the request or suggestion of, a candidate'' (11 C.F.R. \S~109.20(a)). While the Supreme Court previously held in \textit{Colorado Republican Federal Campaign Committee v. FEC} (\textit{Colorado I}), 518 U.S. 604 (1996), that political parties may make unlimited \textit{independent} expenditures without coordinating with candidates, FECA continued to restrict party spending coordinated with candidates. Those statutory caps varied by office and state population: as of 2026, a national party committee could spend between \$130,600 and \$4,071,800 in coordination with an individual Senate candidate and between \$65,300 and \$130,600 in coordination with a House candidate, while the 2024 presidential coordinated-expenditure cap stood at \$32,392,200 (\href{https://www.supremecourt.gov/opinions/25pdf/24-621_h315.pdf}{Slip Opinion, p. 3}).

In 2001, in \textit{FEC v. Colorado Republican Federal Campaign Committee} (\textit{Colorado II}), 533 U.S. 431 (2001), the Supreme Court upheld FECA's party coordinated-expenditure limits by a 5--4 vote, reasoning that they were closely drawn to prevent donors from circumventing individual candidate contribution limits by funneling large donations through political parties. Over the next quarter-century, however, the campaign finance landscape shifted dramatically. In the 2021--2022 election cycle, the National Republican Senatorial Committee (NRSC) spent roughly \$15.5 million and the National Republican Congressional Committee (NRCC) spent roughly \$8.3 million on coordinated expenditures, largely for political advertising (\href{https://www.oyez.org/cases/2025/24-621}{Oyez Case Summary}). Meanwhile, following \textit{Citizens United v. FEC}, 558 U.S. 310 (2010), and \textit{SpeechNow.org v. FEC}, 599 F.3d 686 (D.C. Cir. 2010) (en banc), outside groups such as Super PACs gained the ability to raise and spend unlimited sums on independent expenditures; in the 2023--2024 election cycle, PACs raised over \$15.7 billion compared to \$2.7 billion raised by political parties (\href{https://www.supremecourt.gov/opinions/25pdf/24-621_h315.pdf}{Slip Opinion, p. 25}).

In 2022, two national party committees---the NRSC and the NRCC---together with then-Senate candidate JD Vance and then-Representative Steve Chabot, sued the FEC and its Commissioners, alleging that 52 U.S.C. \S~30116(d) violated the First Amendment both facially and as applied to ``party coordinated communications'' (\href{https://www.supremecourt.gov/opinions/25pdf/24-621_h315.pdf}{Slip Opinion, pp. 3--4}; \href{https://www.oyez.org/cases/2025/24-621}{Oyez Case Summary}). By the time the case reached the Supreme Court, JD Vance had become Vice President of the United States, though he maintained an active Statement of Candidacy with the FEC for a 2028 Senate run and an active principal campaign committee (JD Vance for Senate) that had raised funds for a Senate race (\href{https://www.supremecourt.gov/opinions/25pdf/24-621_h315.pdf}{Slip Opinion, p. 6}).

\section{Procedural History}
Plaintiffs filed suit in 2022 in the U.S. District Court for the Southern District of Ohio, invoking FECA's expedited judicial-review provision, 52 U.S.C. \S~30110, which directs district courts to certify non-frivolous constitutional questions to the regional court of appeals sitting en banc (\href{https://www.brennancenter.org/our-work/research-reports/national-republican-senatorial-committee-v-federal-election-commission}{Brennan Center Case Page}; \href{https://www.oyez.org/cases/2025/24-621}{Oyez Case Summary}). The district court certified the constitutional question in \textit{NRSC v. FEC}, 712 F. Supp. 3d 1017 (S.D. Ohio 2024).

On September 5, 2024, the en banc U.S. Court of Appeals for the Sixth Circuit rejected plaintiffs' facial and as-applied First Amendment challenges, 117 F.4th 389 (6th Cir. 2024) (\href{https://law.justia.com/cases/federal/appellate-courts/ca6/24-3051/24-3051-2024-09-05.html}{Sixth Circuit En Banc Opinion}). Writing for the majority, Chief Judge Sutton held that the lower court remained bound by \textit{Colorado II} unless and until the Supreme Court overruled it, although a majority of Sixth Circuit judges expressed serious doubt about \textit{Colorado II}'s continued vitality in light of \textit{McCutcheon v. FEC}, 572 U.S. 185 (2014), and \textit{FEC v. Ted Cruz for Senate}, 596 U.S. 289 (2022) (\href{https://www.supremecourt.gov/opinions/25pdf/24-621_h315.pdf}{Slip Opinion, p. 4}). Judges Thapar and Bush filed concurrences questioning \textit{Colorado II}; Judges Stranch and Bloomekatz concurred in the judgment defending \textit{Colorado II}; and Judge Readler dissented, arguing that intervening legal and factual changes permitted the Sixth Circuit to invalidate the limits.

Plaintiffs petitioned for certiorari on December 4, 2024 (\href{https://supreme.justia.com/cases/federal/us/609/24-621/}{Justia Case Page}). In May 2025, the federal government declined to defend \S~30116(d), agreeing with plaintiffs that the limits violated the First Amendment. On June 30, 2025, the Supreme Court granted certiorari, 606 U.S. 931 (2025), and granted intervention to the Democratic National Committee (DNC), Democratic Senatorial Campaign Committee (DSCC), and Democratic Congressional Campaign Committee (DCCC); on July 1, 2025, the Court also appointed Roman Martinez as \textit{amicus curiae} to defend the Sixth Circuit's judgment (\href{https://www.supremecourt.gov/opinions/25pdf/24-621_h315.pdf}{Slip Opinion, pp. 4--5}; \href{https://supreme.justia.com/cases/federal/us/609/24-621/}{Justia Case Page}). The case (No. 24-621) was argued on December 9, 2025, and decided on June 30, 2026 (\href{https://www.courtlistener.com/opinion/10882240/national-republican-senatorial-committee-v-federal-election-commn/}{CourtListener Opinion}; \href{https://www.scotusblog.com/cases/national-republican-senatorial-committee-v-federal-election-commission/}{SCOTUSblog Case Page}).

\section{Judicial Votes}
\begin{itemize}
\item \textbf{Vote:} 6--3, reversing and remanding the judgment of the Sixth Circuit (\href{https://www.supremecourt.gov/opinions/25pdf/24-621_h315.pdf}{Slip Opinion, Syllabus, p. 3}).
\item \textbf{Majority (6):} Justice Kavanaugh delivered the opinion of the Court, joined in full by Chief Justice Roberts and Justices Thomas, Alito, Gorsuch, and Barrett.
\item \textbf{Concurrences (0):} None.
\item \textbf{Dissent (3):} Justice Kagan filed a dissenting opinion, joined in full by Justices Sotomayor and Jackson.
\end{itemize}

\section{Holding}
\textbf{Judgment:} Reversed and remanded. The Supreme Court reversed the en banc Sixth Circuit and expressly overruled \textit{FEC v. Colorado Republican Federal Campaign Committee} (\textit{Colorado II}), 533 U.S. 431 (2001) (\href{https://www.supremecourt.gov/opinions/25pdf/24-621_h315.pdf}{Slip Opinion, pp. 26, 29}).

\textbf{Rule Announced:} FECA's limits on political-party coordinated expenditures with candidates, 52 U.S.C. \S~30116(d), violate the First Amendment. Even under ``closely drawn'' scrutiny, a campaign finance restriction must serve the sole permissible governmental interest of preventing \textit{quid pro quo} corruption or its appearance, and it must be necessary, narrowly tailored, and not disproportionate to that interest. Because federal base contribution limits, statutory earmarking rules (52 U.S.C. \S~30116(a)(8)), and modern disclosure requirements (52 U.S.C. \S~30104(b)) already provide three meaningful prophylactic safeguards against circumvention, capping political parties' coordinated expenditures is an unnecessary and disproportionate ``fourth prophylaxis'' that impermissibly abridges core political speech (\href{https://www.supremecourt.gov/opinions/25pdf/24-621_h315.pdf}{Slip Opinion, pp. 9--10, 20--21}). As a threshold jurisdictional holding, the case was not moot despite the Executive Branch's refusal to enforce \S~30116(d) and Vice President Vance's election to executive office, because FECA authorizes private enforcement suits if the FEC fails to act (52 U.S.C. \S\S~30109(a)(1), (a)(8)(A), (a)(8)(C)) and Vice President Vance maintained an active Statement of Candidacy and campaign committee for a 2028 Senate race (\href{https://www.supremecourt.gov/opinions/25pdf/24-621_h315.pdf}{Slip Opinion, pp. 5--6}).

\textbf{Holding vs.\ Dicta:} The Court's binding holdings are: (1) Article III jurisdiction survived because of FECA's private enforcement mechanism and Vice President Vance's active FEC candidacy filings; (2) FECA's coordinated-expenditure caps on political parties in 52 U.S.C. \S~30116(d) fail closely drawn scrutiny and violate the First Amendment; and (3) \textit{Colorado II} is overruled. By contrast, several points are expressly reserved or constitute dicta: (a) the Court explicitly declined to decide whether strict scrutiny or closely drawn scrutiny governs party coordinated expenditures, calling the distinction ``academic'' because the limits fail even closely drawn review (\href{https://www.supremecourt.gov/opinions/25pdf/24-621_h315.pdf}{Slip Opinion, p. 10}); (b) in footnote 6, the Court stated that its decision ``does not address the statutory limits on coordinated expenditures by outside groups'' (\href{https://www.supremecourt.gov/opinions/25pdf/24-621_h315.pdf}{Slip Opinion, p. 26 n. 6}); (c) in footnote 5, the Court noted that it did not rely on Congress's 2014 statutory carveouts for party recount/legal-proceedings funds and convention accounts, though it remarked that those carveouts further undermined the anti-corruption rationale (\href{https://www.supremecourt.gov/opinions/25pdf/24-621_h315.pdf}{Slip Opinion, p. 21 n. 5}); and (d) the Court left intact base contribution limits to candidates and parties, earmarking restrictions, and disclosure obligations (\href{https://www.mayerbrown.com/en/insights/publications/2026/07/supreme-court-strikes-down-limits-on-political-party-coordinated-expenditures}{Mayer Brown Analysis}).

\section{Analysis of Opinions}

\subsection{Majority Opinion (Justice Kavanaugh)}
Justice Kavanaugh's opinion for the six-Justice majority proceeds in five parts (\href{https://www.supremecourt.gov/opinions/25pdf/24-621_h315.pdf}{Slip Opinion, pp. 1--29}; \href{https://www.courtlistener.com/opinion/10882240/national-republican-senatorial-committee-v-federal-election-commn/}{CourtListener Opinion}). First, after summarizing FECA's caps and the procedural posture (Part I), the Court rejected the mootness challenges raised by intervenors and court-appointed \textit{amicus} (Part II), emphasizing both the credible threat of private FECA suits under 52 U.S.C. \S~30109(a)(8) and Vice President Vance's extant 2028 Senate Statement of Candidacy and campaign committee.

Second, on the First Amendment merits (Part III), the Court grounded its analysis in text, history, and modern campaign finance precedent. Because virtually all mass political communication requires spending money, restricting a party's coordinated spending directly abridges its speech and inflicts a ``stifling effect on the ability of the party to do what it exists to do'' (\href{https://www.supremecourt.gov/opinions/25pdf/24-621_h315.pdf}{Slip Opinion, pp. 1, 7--8}). For nearly two centuries after 1791, political parties freely coordinated spending with their nominees without rendering elections corrupt or dysfunctional. Turning to the level of scrutiny, the Court bypassed the choice between strict scrutiny and \textit{Buckley}'s ``closely drawn'' test because FECA's caps fail even closely drawn review, which under \textit{McCutcheon} and \textit{Cruz} demands ``rigorous'' scrutiny showing that a restriction is ``necessary,'' ``narrowly tailored,'' and not ``disproportionate'' (\href{https://www.supremecourt.gov/opinions/25pdf/24-621_h315.pdf}{Slip Opinion, pp. 9--10}).

The majority systematically evaluated four asserted governmental interests: (1) reducing ``wasteful and excessive'' campaign spending (flatly impermissible under \textit{Buckley} and \textit{Cruz}); (2) preventing a party itself from exerting undue influence over its own candidates (nonsensical given parties' inherent role in representative government); (3) curbing wealthy donors' ``undue influence,'' access, or ingratiation (squarely repudiated by \textit{McCutcheon} and \textit{Cruz}, which confine permissible regulation strictly to \textit{quid pro quo} corruption or its appearance); and (4) preventing circumvention of candidate base contribution limits through donations to parties (\href{https://www.supremecourt.gov/opinions/25pdf/24-621_h315.pdf}{Slip Opinion, pp. 10--15}). Focusing on the fourth rationale, the Court reasoned that contributions to a party are one significant step removed from \textit{quid pro quo} corruption because the party controls how funds are spent across many races and priorities. Where a donor actually directs funds to a specific candidate, FECA's earmarking rule (\S~30116(a)(8) and 11 C.F.R. \S~110.6(b)(1)) already counts the donation against the candidate's base limit, while modern internet-era disclosure laws (\S~30104(b)) expose large contributions and coordinated spending to public and regulatory scrutiny. Moreover, a majority of states impose no caps on state-party coordinated expenditures, yet no evidence of resulting corruption has materialized (\href{https://www.supremecourt.gov/opinions/25pdf/24-621_h315.pdf}{Slip Opinion, pp. 16--19}). Adding coordinated-expenditure caps atop base limits, earmarking rules, and disclosure creates a ``fourth prophylaxis'' that is disproportionate and unnecessary.

Third, the majority overruled \textit{Colorado II} under \textit{stare decisis} (Part IV), calling it a ``three-legged stool where all three legs have already been knocked out'' (\href{https://www.supremecourt.gov/opinions/25pdf/24-621_h315.pdf}{Slip Opinion, p. 21}). Applying the \textit{Ramos v. Louisiana} factors, Justice Kavanaugh concluded that: (1) \textit{Colorado II} was egregiously wrong---vindicating Justice Thomas's 2001 dissent---because it used watered-down scrutiny permitting ``unskillful tailoring,'' relied on the now-rejected undue-influence rationale, and mischaracterized parties' unique relationship with candidates; (2) its real-world effects weakened political parties relative to Super PACs and outside groups, fueling fragmentation and polarization; and (3) outside groups had no cognizable reliance interest in preserving an unfair advantage over parties (\href{https://www.supremecourt.gov/opinions/25pdf/24-621_h315.pdf}{Slip Opinion, pp. 22--26}). Finally (Part V), the Court rejected the dissent's joint-fundraising-committee (JFC) concern by invoking \textit{McCutcheon}'s conclusion that any agreement to route JFC funds to a specific candidate triggers existing earmarking rules, and emphasized that the ruling applies equally to all political parties (\href{https://www.supremecourt.gov/opinions/25pdf/24-621_h315.pdf}{Slip Opinion, pp. 26--29}).

\subsection{Dissenting Opinion (Justice Kagan)}
Justice Kagan, joined by Justices Sotomayor and Jackson, argued that the majority dismantled an essential guardrail protecting base candidate contribution limits from massive evasion and disregarded \textit{stare decisis} (\href{https://www.supremecourt.gov/opinions/25pdf/24-621_h315.pdf}{Slip Opinion, Kagan, J., dissenting, pp. 1--19}). Tracing FECA's 1974 amendments to Watergate-era \textit{quid pro quo} scandals such as the dairy industry's \$2 million pledge for milk price supports, Justice Kagan emphasized that both the majority and dissent agree on two core premises: base contribution limits (\$3,500 per election, or \$7,000 for primary and general elections combined) validly prevent \textit{quid pro quo} corruption, and Congress may enact tailored anti-circumvention rules to prevent evasion of those limits (\href{https://www.supremecourt.gov/opinions/25pdf/24-621_h315.pdf}{Slip Opinion, Kagan, J., dissenting, pp. 3--5}).

Justice Kagan challenged the majority's benign portrait of coordinated spending as collaborative ad creation. In practice, she noted, candidates often produce ads on their own and forward the invoice to the party---which replies ``Received. Will process''---or have the party pay pollsters, rent, utilities, and catering bills, making coordinated expenditures ``as useful to the candidate as cash'' and turning the party into ``an alternative checking account for a campaign'' (\href{https://www.supremecourt.gov/opinions/25pdf/24-621_h315.pdf}{Slip Opinion, Kagan, J., dissenting, pp. 2, 6}). Looking ``under the hood of modern fundraising,'' she illustrated how a candidate forms a joint fundraising committee (JFC) with a national party committee (\$44,300 donor cap) and 50 state party committees (\$10,000 cap each), allowing a single wealthy donor to write one check for \$551,300---nearly 80 times the \$7,000 candidate limit---after which state parties immediately transfer their shares to the national party under 52 U.S.C. \S~30116(a)(4) (\href{https://www.supremecourt.gov/opinions/25pdf/24-621_h315.pdf}{Slip Opinion, Kagan, J., dissenting, pp. 7--10}). Citing real-world 2024 FEC data from the Harris Victory Fund and Trump 47 Committee, she demonstrated that this mechanism is no mere hypothetical (\href{https://www.supremecourt.gov/opinions/25pdf/24-621_h315.pdf}{Slip Opinion, Kagan, J., dissenting, p. 12 n. 4}).

Without coordinated-expenditure caps, the national party can spend that entire \$551,300 check directly on the soliciting candidate's bills, and neither earmarking nor disclosure stops the scheme. Earmarking rules do not apply when a donor writes an unearmarked check to a candidate's JFC---even in an explicit \textit{quid pro quo} exchange for a legislative favor---because the donor gives no formal routing instruction, and parties have strong incentives to return funds to prolific candidate-fundraisers (\href{https://www.supremecourt.gov/opinions/25pdf/24-621_h315.pdf}{Slip Opinion, Kagan, J., dissenting, pp. 15--17}). Disclosure, meanwhile, cannot deter corruption that occurs in plain sight, as \textit{Buckley} recognized. Finally, Justice Kagan faulted the majority's \textit{stare decisis} analysis for lacking any ``special justification'': \textit{Colorado II} applied the very same closely drawn standard, rested squarely on \textit{quid pro quo} circumvention, and if parties have weakened relative to Super PACs, the culprit was the Court's own decisions in \textit{Citizens United} and \textit{SpeechNow.org}, not \textit{Colorado II} (\href{https://www.supremecourt.gov/opinions/25pdf/24-621_h315.pdf}{Slip Opinion, Kagan, J., dissenting, pp. 13--15}).

\section{Examples: Future Applications}
\begin{enumerate}
\item \textbf{Same-Side (Invalidated): State-Law Caps on State Party Coordinated Expenditures.} Suppose a state legislature enforces a statute capping how much the state Democratic or Republican party may spend in coordination with its gubernatorial nominee, while also enforcing candidate base contribution limits, anti-earmarking rules, and online campaign finance disclosures. Under \textit{NRSC v. FEC}, this state cap is unconstitutional under the First Amendment (incorporated via the Fourteenth Amendment). Like FECA \S~30116(d), it imposes a disproportionate ``fourth prophylaxis'' on core party-candidate speech when less restrictive tools already combat \textit{quid pro quo} circumvention.

\item \textbf{Same-Side (Invalidated): Federal Cap on Party Coordinated Digital and Field Operations.} Suppose Congress responds to \textit{NRSC v. FEC} by enacting a narrower statute that permits unlimited party coordinated spending on television ads, but caps party coordinated spending on digital micro-targeting, polling, and get-out-the-vote canvassing at \$100,000 per Senate race. This restriction falls on the same side as \textit{NRSC v. FEC} and would be struck down. The majority expressly identified polling, voter targeting, social media, and get-out-the-vote operations as traditional day-to-day party-candidate coordination protected by the First Amendment (\href{https://www.supremecourt.gov/opinions/25pdf/24-621_h315.pdf}{Slip Opinion, p. 8}), and the anti-circumvention rationale fails for the same reasons.

\item \textbf{Opposite-Side (Upheld): Federal Prohibition on Super PAC Coordinated Expenditures.} Suppose a billionaire-funded Super PAC challenges 52 U.S.C. \S~30116(a)(7)(B), arguing that after \textit{NRSC v. FEC}, capping coordinated expenditures by outside independent-expenditure committees likewise violates the First Amendment. This challenge falls on the opposite side and would fail under \textit{NRSC v. FEC}. In footnote 6, the Court explicitly carved out outside groups, emphasizing that political parties possess a unique First Amendment relationship and ``unity of interest'' with their nominees (\href{https://www.supremecourt.gov/opinions/25pdf/24-621_h315.pdf}{Slip Opinion, pp. 24, 26 n. 6}), and outside Super PACs are not subject to the base contribution limits that constrain donors to political parties.

\item \textbf{Opposite-Side (Upheld): FEC Enforcement Action Against Earmarked Party Donations.} Suppose a wealthy donor contributes \$250,000 to a national party committee accompanied by an email instructing the party treasurer that ``this money is strictly to pay for Senator X's October TV buy,'' and the FEC fines the donor and party for violating 52 U.S.C. \S~30116(a)(8) and candidate base contribution limits. This enforcement action falls on the opposite side and is clearly constitutional under \textit{NRSC v. FEC}. Both the majority and petitioners expressly agreed that the government retains a valid interest in enforcing base limits and anti-earmarking rules (\href{https://www.supremecourt.gov/opinions/25pdf/24-621_h315.pdf}{Slip Opinion, pp. 14, 16, 18}).

\item \textbf{Fence-Sitter (Close Case): Federal Cap on Direct Cash Contributions by Political Parties to Candidates (52 U.S.C. \S~30116(a)(2)).} Suppose a national party committee challenges FECA's \$5,000-per-election limit on direct cash transfers from a party to its candidate. On one hand, because \textit{NRSC v. FEC} now allows a party to pay unlimited campaign bills (including rent, catering, and candidate-produced ads) through coordinated expenditures, capping direct cash transfers looks formally anomalous and may be attacked as an unnecessary ``fourth prophylaxis'' given base donor limits, earmarking rules, and disclosure. On the other hand, direct cash transfers have historically been governed by \textit{Buckley}'s deferential contribution-limit framework, involve zero expressive content of the party's own, and pose an even more direct conduit risk than coordinated speech---making this a genuine doctrinal fence-sitter after \textit{NRSC v. FEC}.
\end{enumerate}

\section{Critique}
\textit{NRSC v. FEC} generated an unusually cross-cutting scholarly response, dividing traditional ideological camps over the constitutional role of political parties. From an \textbf{originalist and textualist perspective}, supporters praised Justice Kavanaugh's fidelity to the First Amendment's text (``Congress shall make no law \ldots\ abridging the freedom of speech'') and the nearly 200-year historical tradition of unrestricted party-candidate coordination prior to 1974 (\href{https://www.supremecourt.gov/opinions/25pdf/24-621_h315.pdf}{Slip Opinion, pp. 8, 28}). Defenders of robust political-speech rights argued that \textit{Colorado II} had always been an anomaly after \textit{Citizens United}, \textit{McCutcheon}, and \textit{Cruz}, and that overruling it eliminated a distorted regime that favored unaccountable Super PACs over transparent, broad-based political parties (\href{https://www.mayerbrown.com/en/insights/publications/2026/07/supreme-court-strikes-down-limits-on-political-party-coordinated-expenditures}{Mayer Brown Analysis}). Notably, even several centrist and reform-minded election-law scholars---including Samuel Issacharoff, Richard Pildes, and Bob Bauer---welcomed the outcome on institutional grounds. Issacharoff called the ruling ``a welcome and overdue correction'' that ended an ``artificial constitutional distinction between parties and the candidates they exist to elect'' and began rebalancing the ``hydraulic system'' of campaign finance away from unaccountable outside groups toward governing party institutions (\href{https://democracyproject.org/posts/a-welcome-correction-nrsc-v-fec-and-the-need-to-empower-parties}{Issacharoff, Democracy Project}; \href{https://electionlawblog.org/2026/issacharoff-a-welcome-correction-nrsc-v-fec-and-the-need-to-empower-parties/}{Election Law Blog}).

Conversely, \textbf{progressive and doctrinal critics}---including the Brennan Center for Justice and the Constitutional Accountability Center---attacked both the Court's institutional overreach and its empirical assumptions (\href{https://www.brennancenter.org/our-work/research-reports/national-republican-senatorial-committee-v-federal-election-commission}{Brennan Center Case Page}; \href{https://www.theusconstitution.org/litigation/national-republican-senatorial-committee-v-federal-election-commission/}{Constitutional Accountability Center}). Even while acknowledging that strengthening political parties might be defensible as legislative policy, these critics argued that Congress---not an unelected Court---should make complex structural trade-offs in campaign finance regulation, and that the majority once again brushed aside \textit{stare decisis} without a genuine special justification. Substantively, critics echoed Justice Kagan's warning that the majority ignored the real-world mechanics of modern joint fundraising committees: by combining \textit{McCutcheon}'s abolition of aggregate limits with unlimited party coordinated spending, the Court effectively raised the functional ceiling on a single donor's candidate support from \$7,000 to over \$550,000, leaving earmarking and disclosure laws unable to police tacit \textit{quid pro quo} arrangements (\href{https://www.theusconstitution.org/litigation/national-republican-senatorial-committee-v-federal-election-commission/}{Constitutional Accountability Center}).

At bottom, the case exposes a deep clash of \textbf{competing constitutional values}: (1) \textit{robust First Amendment protection for associative political speech} and the institutional vitality of political parties versus (2) \textit{democratic integrity, anti-corruption prophylaxis, and legislative deference} to Congress's judgment about how money actually flows in modern elections.

\section{Key Quotations}

\begin{quote}
``The base limits on contributions to candidates serve as an initial prophylaxis against quid pro quo corruption or its appearance in this context---after all, most contributions to candidates are not given in exchange for some official action. \textit{Id.}, at 221. The earmarking rules constitute a second prophylaxis. The disclosure requirements supply a third prophylaxis. So prophylaxis upon prophylaxis upon prophylaxis already serve to prevent quid pro quo corruption or its appearance.''
\end{quote}
\noindent --- Justice Kavanaugh, Opinion of the Court, \href{https://www.supremecourt.gov/opinions/25pdf/24-621_h315.pdf}{609 U.S. \_\_\_, slip op. at 20 (2026)} (also on \href{https://www.courtlistener.com/opinion/10882240/national-republican-senatorial-committee-v-federal-election-commn/}{CourtListener}).

\begin{quote}
``Courts cannot simply say, `what's the harm in allowing just one more regulation' when that regulation would limit freedom of speech. On the contrary, courts must preserve and protect the freedom of speech guaranteed by the Framers.''
\end{quote}
\noindent --- Justice Kavanaugh, Opinion of the Court, \href{https://www.supremecourt.gov/opinions/25pdf/24-621_h315.pdf}{609 U.S. \_\_\_, slip op. at 20 (2026)}.

\begin{quote}
``To sum up: In light of the other meaningful prophylactic measures available to the Government, and given the severe infringement on First Amendment-protected political speech that ensues from limiting a political party's spending in support of its candidates, we conclude that the political-party coordinated-expenditure limits are `disproportionate' and are not `necessary' and `narrowly tailored' for the circumvention interest it seeks to protect.''
\end{quote}
\noindent --- Justice Kavanaugh, Opinion of the Court, \href{https://www.supremecourt.gov/opinions/25pdf/24-621_h315.pdf}{609 U.S. \_\_\_, slip op. at 21 (2026)}.

\begin{quote}
``The majority invalidates Congress's restriction of coordinated expenditures, thus enabling a party to serve as an alternative checking account for a campaign. As a result, a donor will be able to give a party as much as half a million dollars (as compared to the \$7,000 he can give directly to the candidate) to cover the candidate's bills. And the candidate can seek just such a donation. So the Court ushers back in the same opportunities for quid pro quo corruption that the contribution limits were meant to check.''
\end{quote}
\noindent --- Justice Kagan, Dissenting Opinion, \href{https://www.supremecourt.gov/opinions/25pdf/24-621_h315.pdf}{609 U.S. \_\_\_, slip op. at 2 (2026)}.

\begin{quote}
``When this Court in \textit{McCutcheon} invalidated aggregate limits, Justice Breyer wrote in dissent: `[T]oday's decision eviscerates our Nation's campaign finance laws, leaving a remnant incapable of dealing with the grave problems of democratic legitimacy that those laws were intended to resolve.' 572 U.S., at 233. I'm not sure what to call a remnant of a remnant, but that is what the Court has left today.''
\end{quote}
\noindent --- Justice Kagan, Dissenting Opinion, \href{https://www.supremecourt.gov/opinions/25pdf/24-621_h315.pdf}{609 U.S. \_\_\_, slip op. at 19 (2026)}.

\end{document}
